\section{方法}

\subsection{图的定义}

令人物图为 $G=(V,E)$：顶点 $V$ 是人物表中的实体，共 1231 个；
边 $e\in E$ 是关系表中的一条记录，端点为关系两端的裸人物名，共 2324 条。
边无向化：关系类别本身带方向语义（如「君臣从属」有上下之分），
但本篇只用量类别构成与连通性，方向不参与任何计数。
自环与端点不在 $V$ 内的记录剔除。

\subsection{集中度：基尼系数}

对任意度数列 $d$（按升序排列），基尼系数按
\begin{equation}\label{eq:gini}
\mathrm{Gini}(d)=\frac{2\sum_{i=1}^{n} i\,d_{(i)}}{n\sum_{i=1}^{n}d_{(i)}}-\frac{n+1}{n}
\end{equation}
计算。本篇报两个口径：全图 1231 人的基尼（含孤立点），
以及仅取有边顶点的子列的基尼——两者不相等，
孤立点会把全图基尼往高里推，把它们并进来比较会得出"比看上去更集中"的错误结论。

\subsection{分布形状：Hill 尾指数与它的漂移}

对截断 $x_{\min}$，取所有 $d\ge x_{\min}$ 的观测，Hill 估计器为
\begin{equation}\label{eq:hill}
\hat\alpha = 1+\left[\sum_{d\ge x_{\min}}\ln\frac{d}{x_{\min}}\right]^{-1}.
\end{equation}
按 \citet{clauset2009} 的口径，幂律主张者必须证明 $\hat\alpha$ 在多个 $x_{\min}$
上稳定、且尾样本足够大；本篇因此不取"某一个最优 xmin 的 $\hat\alpha$"作结论，
而报 xmin = 2, 3, 5, 8, 10, 15 六个切点上的 $\hat\alpha$ 与尾样本量，
并以极差（最大值减最小值）作为漂移量。漂移量与尾样本量一起决定能否贴标签。

\subsection{连通块}

在有向边无向化后的图上做一次连通分量分解，报分量个数、最大分量规模，
以及最大分量占（a）带边顶点、（b）全部人物两者的比例。
这两个比例常常差很多，正是本篇要报的一处反差。

\subsection{时间分布与时期迁移}

事件年份序列的报告量为：带年份事件数、覆盖年份数、年均 $\mu$、
样本标准差 $s$、变异系数 $s/\mu$、以及按年事件数的基尼。
$10$ 年为一组作十年汇总供作图（图 \ref{fig:ts} 用的是十年桶）。

时期切分按王朝中段固定为三刀：P1 = 1368--1499、P2 = 1500--1599、P3 = 1600--1662。
这一刀不来自数据聚类，而是先验的王朝分期；切分固定之后，
对「时期 $\times$ 事件类别」列联表做Pearson 卡方独立性检验：
\begin{equation}\label{eq:chi2}
\chi^{2}=\sum_{c,\,p}\frac{(O_{cp}-E_{cp})^{2}}{E_{cp}},
\qquad
E_{cp}=\frac{n_{c}\cdot n_{p}}{n},
\end{equation}
其中 $n_c$ 为类别 $c$ 在三期的合计、$n_p$ 为时期 $p$ 的事件数、$n$ 为总事件数；
自由度 $(3-1)(C-1)$，$C$ 为参与计算的类别数。
效应量报 Cramér's V，
\begin{equation}\label{eq:v}
V=\sqrt{\frac{\chi^{2}}{n\,(\min(3,\,C)-1)}}.
\end{equation}
参与计算的类别限为三期中有任一期出现过的类别；
只在一期出现的类别会独占卡方贡献，本篇因此把单期独类的行一并列出供核对。

\subsection{可复算性}

本节全部量由 \texttt{.dump/}\_stats\_systems.json 提供，正文出现的每个小数
都能在该文件中找到；作图脚本独立存于 figures/，重跑即可复现图 1--3。
本篇不跑任何原书全文检索，因此不受原书不外发的限制；
但事件类别标签与人物抽取由前序工序产出，本篇不声称独立重建了它们。
